%% ma: L11-B14-0030
%% lop: 11
%% chuong: 4
%% bai: 14
%% dang: 11.14.4
%% loai: TLN
%% muc: VD
%% dap_an: "0,25"
%% nguon: "tai-lieu-hoc-tap-mon-toan-lop-11-theo-chuong-trinh-giao-duc-pho-thong-2018.pdf (Lớp Toán Thầy Duy, Phạm Lê Duy, THPT Chu Văn An), Chương 4 Bài 5, Phần 3, Câu 6 (tr.535)"
%% kien_thuc: "Lớp 11 Bài 14 (phép chiếu song song)"
%% trang_thai: cho-duyet
%% ly_do: "Dạng mới đề xuất (11.14.x), chờ giáo viên duyệt dạng; tài liệu gốc không có lời giải mẫu, đáp án do kho tự giải và kiểm bằng sympy."
\begin{ex}
Cho hình chóp $S.ABCD$ có đáy là hình bình hành, trên cạnh $SA$ lấy điểm $M$ sao cho $MA=2MS$. Gọi $O$ là tâm của hình bình hành $ABCD$. Một phép chiếu song song theo phương $MO$ lên mặt phẳng $(ABCD)$ biến điểm $S$ thành điểm $N$. Tỉ số $\dfrac{CN}{CA}$ bằng
\shortans{0,25}
\loigiai{\sloppy $N$ là giao của $AC$ với đường thẳng qua $S$ song song $MO$ (trong mặt phẳng $(SAC)$). Trong $\triangle ASN$, $MO\parallel SN$ với $M\in AS$, $O\in AN$ nên $\dfrac{AM}{AS}=\dfrac{AO}{AN}$. Do $AM=\dfrac23AS$ nên $AN=\dfrac32AO=\dfrac34AC$, suy ra $CN=AC-AN=\dfrac14AC$ và $\dfrac{CN}{CA}=0{,}25$.}
\end{ex}
